\documentclass[11pt]{article}
\usepackage[T1]{fontenc}
\usepackage[margin=1in]{geometry}
\usepackage{amsmath,amssymb,amsthm,microtype}
\usepackage[hidelinks]{hyperref}
\newtheorem{theorem}{Theorem}
\newcommand{\ellp}{\ell_p^{,2}}
\newcommand{\gap}{\Gamma_p}
\title{A non-virtually-degenerate $p$-polygonal equality in $\ell_p^2$ for every $p>2$}
\author{Autonomous counterexample for arXiv:1203.5837}
\date{11 August 2026}

\begin{document}
\maketitle

\begin{abstract}
Kelleher--Miller--Osborn--Weston proved that virtual degeneracy is sufficient
for a nontrivial $p$-polygonal equality in every $L_p$ space and asked whether
it is necessary when $p>2$.  It is not.  For every $p>2$ we give a
three-vertex weighted simplex in $\ell_p^2$ whose two scalar-coordinate gaps
have opposite signs and cancel after one explicit rescaling.  One coordinate
simplex is nondegenerate, so the equality is not virtually degenerate.
\end{abstract}

\section{Gap convention and result}

For a normalized signed simplex
\[
 D=[x_i(m_i);y_j(n_j)]_{s,t},\qquad
 m_i,n_j>0,\quad \sum_i m_i=\sum_j n_j=1,
\]
put
\begin{align*}
 \gap(D)={}&\sum_{i,j}m_in_j\|x_i-y_j\|_p^p
 -\sum_{i<i'}m_im_{i'}\|x_i-x_{i'}\|_p^p\\
 &-\sum_{j<j'}n_jn_{j'}\|y_j-y_{j'}\|_p^p.
\end{align*}
A nondegenerate simplex with $\gap(D)=0$ is a nontrivial
$p$-polygonal equality in the terminology of the source.

\begin{theorem}\label{thm:main}
For every $p>2$, virtual degeneracy is not necessary for a nontrivial
$p$-polygonal equality in an $L_p$ space.  More precisely, $\ell_p^2$ contains
a completely refined signed $(2,1)$-simplex $D$ such that $\gap(D)=0$ but $D$
is not virtually degenerate.
\end{theorem}

\section{Construction and proof}

Fix $p>2$ and define $c=c_p>0$ by
\begin{equation}\label{eq:c}
 c^p=\frac{2^{p-2}-1}{2^{p-1}+\tfrac14}.
\end{equation}
The numerator is positive precisely because $p>2$.  In $\ellp$ set
\begin{equation}\label{eq:vertices}
 x_1=(-1,0),\qquad x_2=(1,c),\qquad y_1=(0,2c),
 \qquad m_1=m_2=\frac12,\quad n_1=1.
\end{equation}
All three vertices are distinct and all weights are positive, so
$D=[x_1(1/2),x_2(1/2);y_1(1)]$ is completely refined and in particular
nondegenerate.

There is only one within-side term.  Directly from \eqref{eq:vertices},
\begin{align*}
 \gap(D)
 &=\frac12\bigl(1+(2c)^p\bigr)
   +\frac12\bigl(1+c^p\bigr)
   -\frac14\bigl(2^p+c^p\bigr)\\
 &=1-2^{p-2}+\left(2^{p-1}+\frac14\right)c^p=0,
\end{align*}
where the last equality is \eqref{eq:c}.  Thus $D$ gives a nontrivial
$p$-polygonal equality.

It remains only to exclude virtual degeneracy.  In the first scalar
coordinate the induced signed simplex is
\[
 [-1(1/2),,1(1/2);,0(1)].
\]
Its three vertices are distinct.  At the scalar point $-1$, for example, the
total weight on the $x$-side is $1/2$ and that on the $y$-side is zero.  Hence
this scalar simplex is nondegenerate.  Virtual degeneracy in $\ell_p^2$
requires degeneracy in every scalar coordinate, so $D$ is not virtually
degenerate.  This proves Theorem~\ref{thm:main}. \qed

\section{Mechanism and scope}

The construction is a cancellation that is unavailable below exponent two.
The first coordinate contributes
\[
 1-2^{p-2}<0
\]
to the gap.  Before scaling, the second coordinate---the scalar simplex
$[0(1/2),1(1/2);2(1)]$---contributes
\[
 2^{p-1}+\frac14>0.
\]
Because $\|z\|_p^p$ is additive over coordinates and scaling a coordinate by
$c$ multiplies its gap by $c^p$, equation \eqref{eq:c} cancels the two
contributions exactly.

This answers only the explicit necessity question in the first paragraph of
the source's open-problems section (source PDF page 17).  It does not classify
all $p$-polygonal equalities for $p>2$.  The later arXiv:2404.06658 proves a
general existence criterion and settles the source's separate Schatten-class
existence question, but does not discuss virtual degeneracy; the elementary
counterexample above is independent of that result.

\end{document}
